\textbf{Phần 1. Trắc nghiệm lựa chọn}
\setcounterpageref{ex}{0}

\begin{ex} %[2D13V65]
Trên biểu đồ tranh, mỗi biểu tượng ngôi sao biểu thị cho $2$ điểm tốt. Bạn Mai có $5$ biểu tượng ngôi sao. Mai có bao nhiêu điểm tốt?
\choice
{$7$ điểm}
{$5$ điểm}
{\True $10$ điểm}
{$12$ điểm}
\loigiai{
$2 \times 5 = 10$ (điểm tốt).
\textbf{Đáp án của câu hỏi là C}
}
\end{ex}

\begin{ex} %[2D13V65]
Một biểu đồ tranh biểu diễn số cây trồng của các tổ, mỗi hình cái cây đại diện cho $5$ cây xanh. Tổ Một trồng được $4$ hình cái cây. Tổ Một trồng được bao nhiêu cây xanh?
\choice
{\True $20$ cây}
{$9$ cây}
{$15$ cây}
{$25$ cây}
\loigiai{
$5 \times 4 = 20$ (cây xanh).
\textbf{Đáp án của câu hỏi là A}
}
\end{ex}

\begin{ex} %[2D13V65]
Biểu đồ tranh hoa điểm mười của 4 bạn: An có 6 bông, Bình có 4 bông, Cúc có 8 bông, Dũng có 5 bông (mỗi bông biểu thị 1 điểm 10). Bạn đạt nhiều điểm 10 nhất là
\choice
{An}
{Bình}
{Dũng}
{\True Cúc}
\loigiai{
Cúc có 8 bông hoa điểm 10, nhiều nhất.
\textbf{Đáp án của câu hỏi là D}
}
\end{ex}

\begin{ex} %[2D13V65]
Trên biểu đồ tranh bán xe đạp, mỗi hình xe đạp biểu thị $10$ chiếc xe. Cửa hàng vẽ $6$ hình xe đạp. Cửa hàng đã bán được bao nhiêu chiếc xe đạp?
\choice
{$16$ chiếc}
{\True $60$ chiếc}
{$50$ chiếc}
{$6$ chiếc}
\loigiai{
$10 \times 6 = 60$ (chiếc xe đạp).
\textbf{Đáp án của câu hỏi là B}
}
\end{ex}

\begin{ex} %[2D13V65]
Biểu đồ tranh số con cá câu được của ba bạn: Nam câu được 3 con, Việt câu được 5 con, Robot câu được 4 con. Cả ba bạn câu được bao nhiêu con cá?
\choice
{$10$ con}
{$11$ con}
{\True $12$ con}
{$14$ con}
\loigiai{
$3 + 5 + 4 = 12$ (con cá).
\textbf{Đáp án của câu hỏi là C}
}
\end{ex}

\begin{ex} %[2D13V65]
Trên biểu đồ số bạn đi học bằng các phương tiện: Đi bộ có 8 bạn, đi xe đạp có 12 bạn, đi ô tô buýt có 15 bạn. Số bạn đi bộ ít hơn số bạn đi xe đạp là
\choice
{\True $4$ bạn}
{$3$ bạn}
{$5$ bạn}
{$7$ bạn}
\loigiai{
$12 - 8 = 4$ (bạn).
\textbf{Đáp án của câu hỏi là A}
}
\end{ex}

\begin{ex} %[2D13V65]
Mỗi biểu tượng quả táo biểu thị $2\text{ kg}$ táo. Nếu muốn biểu thị $16\text{ kg}$ táo thì cần vẽ bao nhiêu biểu tượng quả táo?
\choice
{$6$ biểu tượng}
{$10$ biểu tượng}
{$14$ biểu tượng}
{\True $8$ biểu tượng}
\loigiai{
$16 : 2 = 8$ (biểu tượng).
\textbf{Đáp án của câu hỏi là D}
}
\end{ex}

\begin{ex} %[2D13V65]
Biểu đồ số con vật nuôi trong trang trại: Bò có 20 con, dê có 15 con, cừu có 25 con. Tổng số con vật nuôi là
\choice
{$50$ con}
{\True $60$ con}
{$55$ con}
{$65$ con}
\loigiai{
$20 + 15 + 25 = 60$ (con).
\textbf{Đáp án của câu hỏi là B}
}
\end{ex}

\begin{ex} %[2D13V65]
Trên biểu đồ tranh, mỗi khuôn mặt cười biểu thị cho $5$ bạn học sinh yêu thích môn Toán. Lớp có $7$ khuôn mặt cười. Có bao nhiêu bạn yêu thích môn Toán?
\choice
{$30$ bạn}
{$40$ bạn}
{\True $35$ bạn}
{$12$ bạn}
\loigiai{
$5 \times 7 = 35$ (bạn).
\textbf{Đáp án của câu hỏi là C}
}
\end{ex}

\begin{ex} %[2D13V65]
Biểu đồ số sách đọc trong tuần: Tuần 1 đọc 4 cuốn, tuần 2 đọc 6 cuốn, tuần 3 đọc 5 cuốn, tuần 4 đọc 7 cuốn. Trung bình hoặc tổng số sách đọc trong 4 tuần là
\choice
{\True $22$ cuốn}
{$20$ cuốn}
{$24$ cuốn}
{$18$ cuốn}
\loigiai{
Tổng số sách: $4 + 6 + 5 + 7 = 22$ (cuốn).
\textbf{Đáp án của câu hỏi là A}
}
\end{ex}

\begin{ex} %[2D13V65]
Biểu đồ số ly kem bán được trong ngày: Buổi sáng có 6 hình ly kem, buổi chiều có 9 hình ly kem (mỗi hình ly kem ứng với 2 ly thật). Buổi chiều bán nhiều hơn buổi sáng bao nhiêu ly kem?
\choice
{$3$ ly}
{$4$ ly}
{$8$ ly}
{\True $6$ ly}
\loigiai{
Nhiều hơn: $(9 - 6) \times 2 = 6$ ly kem.
\textbf{Đáp án của câu hỏi là D}
}
\end{ex}

\begin{ex} %[2D13V65]
Khi đọc biểu đồ tranh, thông tin nào sau đây là quan trọng nhất cần xem trước tiên để hiểu tỉ lệ?
\choice
{Tên người vẽ biểu đồ}
{\True Mỗi biểu tượng biểu thị cho bao nhiêu đối tượng}
{Màu sắc của biểu tượng}
{Kích thước của khung tranh}
\loigiai{
Cần xem quy ước mỗi biểu tượng biểu thị cho bao nhiêu đơn vị để tính đúng.
\textbf{Đáp án của câu hỏi là B}
}
\end{ex}

\textbf{Phần 2. Trắc nghiệm đúng sai}
\setcounterpageref{ex}{0}

\begin{ex} %[2D13V65]
Xét các nhận định về đọc và biểu diễn biểu đồ tranh:
\choiceTF
{\True Biểu đồ tranh dùng các hình ảnh hoặc biểu tượng để biểu diễn số liệu.}
{\True Một biểu tượng trên biểu đồ tranh có thể đại diện cho 1, 2, 5 hoặc 10 đối tượng tùy theo quy ước.}
{Biểu đồ tranh không thể dùng để so sánh số lượng giữa các nhóm.}
{\True Muốn biết số lượng thực tế, ta nhân số biểu tượng với giá trị của mỗi biểu tượng.}
\loigiai{
\begin{itemize}
    \item Ý a) Đúng: Đúng định nghĩa biểu đồ tranh.
    \item Ý b) Đúng: Đúng quy ước biểu diễn.
    \item Ý c) Sai: Biểu đồ tranh giúp so sánh rất trực quan. Sai.
    \item Ý d) Đúng: Đúng phương pháp tính.
\end{itemize}
\begin{itemize}
    \item a) \textbf{Đúng}.
    \item b) \textbf{Đúng}.
    \item c) \textbf{Sai}.
    \item d) \textbf{Đúng}.
\end{itemize}
}
\end{ex}

\begin{ex} %[2D13V65]
Biểu đồ tranh số quả cam thu hoạch được: Cây 1 có 3 hình sọt cam, cây 2 có 5 hình sọt cam. Mỗi hình sọt cam biểu thị 10 quả cam:
\choiceTF
{\True Cây 1 thu hoạch được 30 quả cam.}
{\True Cây 2 thu hoạch được 50 quả cam.}
{Cả hai cây thu hoạch được 70 quả cam.}
{\True Cây 2 thu hoạch nhiều hơn cây 1 là 20 quả cam.}
\loigiai{
\begin{itemize}
    \item Ý a) Đúng: 3 x 10 = 30 quả. Đúng.
    \item Ý b) Đúng: 5 x 10 = 50 quả. Đúng.
    \item Ý c) Sai: 30 + 50 = 80 quả, không phải 70 quả. Sai.
    \item Ý d) Đúng: 50 - 30 = 20 quả. Đúng.
\end{itemize}
\begin{itemize}
    \item a) \textbf{Đúng}.
    \item b) \textbf{Đúng}.
    \item c) \textbf{Sai}.
    \item d) \textbf{Đúng}.
\end{itemize}
}
\end{ex}

\textbf{Phần 3. Trả lời ngắn}
\setcounterpageref{ex}{0}

\begin{ex} %[2D13V65]
Trên biểu đồ tranh, mỗi hình bông hoa biểu thị 5 điểm mười. Tổ Một có 6 hình bông hoa. Tổ Một có tất cả bao nhiêu điểm mười?
\shortans{30}
\loigiai{
\textbf{Đáp số:} $30$\\
$5 \times 6 = 30$ điểm mười.
}
\end{ex}

\begin{ex} %[2D13V65]
Để biểu diễn 40 chiếc xe ô tô trên biểu đồ tranh, biết mỗi hình chiếc xe biểu thị 10 chiếc. Cần vẽ bao nhiêu hình chiếc xe?
\shortans{4}
\loigiai{
\textbf{Đáp số:} $4$\\
$40 : 10 = 4$ hình.
}
\end{ex}

\begin{ex} %[2D13V65]
Biểu đồ số bóng bàn bán được: Ngày thứ nhất có 8 quả, ngày thứ hai có 12 quả (mỗi hình vẽ 1 quả bóng). Cả hai ngày bán được bao nhiêu quả?
\shortans{20}
\loigiai{
\textbf{Đáp số:} $20$\\
$8 + 12 = 20$ quả.
}
\end{ex}

\begin{ex} %[2D13V65]
Mỗi biểu tượng cuốn sách đại diện cho 2 cuốn sách. Mai vẽ 9 biểu tượng cuốn sách. Mai có bao nhiêu cuốn sách?
\shortans{18}
\loigiai{
\textbf{Đáp số:} $18$\\
$2 \times 9 = 18$ cuốn.
}
\end{ex}

\begin{ex} %[2D13V65]
Tổ Hai trồng được 25 cây xanh, trên biểu đồ mỗi biểu tượng đại diện cho 5 cây. Cần vẽ mấy biểu tượng?
\shortans{5}
\loigiai{
\textbf{Đáp số:} $5$\\
$25 : 5 = 5$ biểu tượng.
}
\end{ex}

\begin{ex} %[2D13V65]
Trên biểu đồ tranh, bạn Hải có 7 ngôi sao, bạn Nam có 4 ngôi sao (mỗi ngôi sao là 2 điểm tốt). Hải nhiều hơn Nam bao nhiêu điểm tốt?
\shortans{6}
\loigiai{
\textbf{Đáp số:} $6$\\
$(7 - 4) \times 2 = 6$ điểm tốt.
}
\end{ex}
